\section{Dataset and Preprocessing}

The experimental corpus consisted of five PDF documents related to biomedical natural language processing, transformer-based clinical language models, electronic health record analysis, ICD code prediction, BioBERT, ClinicalBERT, and large language models in healthcare. These documents were selected because they contained technical content relevant to evaluating retrieval-augmented generation for biomedical and healthcare-related question answering.

The source PDF files were uploaded to an Amazon S3 bucket and synchronized with an Amazon Bedrock Knowledge Base. During synchronization, the documents were parsed using the default Bedrock parser, segmented into fixed-size chunks of 300 tokens with a 20$\%$ overlap, embedded using Amazon Titan Text Embeddings V2, and indexed in Amazon OpenSearch Serverless. No manual text extraction or external preprocessing was applied before ingestion, allowing the document-processing workflow to remain consistent with the managed Bedrock Knowledge Base configuration.

The evaluation query set consisted of 25 natural-language questions covering biomedical NLP, transformer models, clinical text classification, ICD coding, BioBERT, ClinicalBERT, and healthcare-focused language models. Each query was processed independently through the same hybrid retrieval, reranking, answer-generation, and grounding-evaluation pipeline.

\subsection{Justification for Pilot-Scale Evaluation}
\label{subsec:pilot_scale_justification}
This study was designed as a pilot-scale proof-of-concept evaluation using five source documents and 25 evaluation queries. The purpose of this experimental design was to examine the feasibility, traceability, and grounding behavior of the proposed hybrid retrieval and reranking framework before expanding to a larger benchmark-scale evaluation. The selected documents covered the main technical areas required for the study, including biomedical natural language processing, transformer-based clinical language models, ICD code prediction, BioBERT, ClinicalBERT, and large language models in healthcare. Using a limited but focused document collection allowed the retrieval, reranking, answer-generation, and grounding-evaluation stages to be examined systematically while maintaining interpretability of the retrieved evidence. The 25-query set was designed to provide sufficient topical variation for evaluating the proposed pipeline without making the experiment unnecessarily large during the initial validation stage.

%This study was designed as a pilot-scale proof-of-concept evaluation using five source documents and 25 evaluation queries. The purpose of this experimental design was to examine the feasibility, traceability, and grounding behavior of the proposed hybrid retrieval and reranking framework before expanding to a larger benchmark-scale evaluation. This study used five source documents and 25 evaluation queries as a controlled pilot-scale experimental design. The selected documents covered the main technical areas required for the study, including biomedical natural language processing, transformer-based clinical language models, ICD code prediction, BioBERT, ClinicalBERT, and large language models in healthcare. Using a limited but focused document collection allowed the retrieval, reranking, answer-generation, and grounding-evaluation stages to be examined systematically while maintaining interpretability of the retrieved evidence. The 25-query set was designed to provide sufficient topical variation for evaluating the proposed pipeline without making the experiment unnecessarily large during the initial validation stage.

The decision to use a pilot-scale dataset was also influenced by the cost structure of managed RAG services. Amazon Bedrock pricing is usage-based and depends on model inference, input tokens, output tokens, embedding generation, and other service components used during the workflow \cite{awsbedrockpricing}. In addition, Amazon OpenSearch Serverless charges separately for compute and storage resources, with compute capacity measured in OpenSearch Compute Units \cite{opensearchpricing}. Because each experimental query required hybrid retrieval, reranking, answer generation, and claim-level evaluation, increasing the number of documents or queries would increase the number of API calls, retrieved context tokens, generated tokens, and judge-model evaluations. Therefore, the five-document and 25-query configuration provided a practical balance between experimental coverage, reproducibility, and cost control for an initial proof-of-concept evaluation. Future work will expand the number of documents and queries after the pipeline is further optimized and baseline comparisons are added.
